\section{Traffic Analysis}\label{sec:traffic}

The required course module is Traffic Analysis, delivered as PyShark analyzer
code that exits non-zero when a wire property fails, not as screenshots.
% src: facts-F6#1, facts-F6#2, facts-F6#3, facts-F6#32
Capture runs only on allowlisted loopback or Docker-bridge interfaces.
% src: facts-F6#9, facts-F6#10, facts-F6#12, facts-F6#13
One scenario is captured over plain HTTP and again under TLS, the Postgres
link included.
% src: facts-F6#18, facts-F6#20
The analyzer reassembles TCP streams and searches 13 canaries (record
plaintext, keys, shares) as raw bytes, hex, base64 and PEM; a lost segment
fails the check.
% src: facts-F6#24, facts-F6#25, facts-F6#27, facts-F6#28
The captures come from a slightly earlier version of the code than the
tests (Section~\ref{sec:eval}, Table~\ref{tab:wire}).
% src: facts-F6#34, facts-F6#35, facts-F6#36, facts-F6#37, facts-F6#38, facts-F6#39, facts-F6#40, facts-F6#41, facts-F6#42, facts-F6#43, facts-F6#44, facts-F6#45, facts-F6#46, facts-F6#47, facts-F6#48, facts-F6#49, facts-F6#50, facts-F6#51, facts-F6#52, facts-F6#53, facts-F6#90

% Generated by paper/data/make_wire.py from bench/results/p55_wire_checks.json.
% Do not edit by hand.
\begin{table}[!t]
\centering
\footnotesize
\setlength{\tabcolsep}{3pt}
\caption{PyShark wire checks on the plain and TLS captures (provenance in Section~\ref{sec:eval}).}
% src: facts-F6#34, facts-F6#35, facts-F6#36, facts-F6#37, facts-F6#38, facts-F6#39, facts-F6#40, facts-F6#41, facts-F6#42, facts-F6#43, facts-F6#44, facts-F6#45, facts-F6#46, facts-F6#47, facts-F6#48, facts-F6#49, facts-F6#50, facts-F6#51, facts-F6#52, facts-F6#53, facts-F6#90, facts-F6#93
\label{tab:wire}
\begin{tabular}{@{}l >{\raggedright\arraybackslash}p{0.19\columnwidth} >{\raggedright\arraybackslash}p{0.29\columnwidth} >{\raggedright\arraybackslash}p{0.27\columnwidth} l@{}}
\toprule
\# & Property & Plain capture & TLS capture & Result \\
\midrule
0 & Preconditions & 16 endpoints seen & TLS records only & pass \\
1 & No plaintext & 6 of 13 canaries hit (control fired) & 0 of 13 canaries hit & pass \\
2 & Commit before release & COMMIT 1.35\,ms before response & n/a (plain only) & pass \\
3 & Search leakage & 1 search, 1 term in clear, 43\,B resp. & n/a (plain only) & measured$^\ast$ \\
4 & Break-glass as metadata & URI visible; 4314\,B request & 4356\,B burst; 0 content hits & pass \\
\bottomrule
\multicolumn{5}{@{}p{\columnwidth}@{}}{$^\ast$Leakage is measured, not asserted; the check only confirms search traffic was captured.}\\
\end{tabular}
\end{table}


\emph{No plaintext on the wire holds under TLS only.}
% src: facts-F6#6, facts-F6#35
No canary appears under TLS; the plain capture is the positive control,
exposing record plaintext, the emergency key and four attribute keys, and a
clean TLS result is not trusted unless that control fires.
% src: facts-F6#35, facts-F6#36, facts-F6#37
Even on the plain wire no share, signing key or DEK appears.
% src: facts-F6#6, facts-F6#37

\emph{Commit before release.} The analyzer matches the audit service's
Postgres INSERT and \texttt{COMMIT} acknowledgement to the
share-release response, inside which the sealed audit shares
must first appear (the analyzer keys on share~3; share~4 travels in the same
response and is not timed separately); the \texttt{COMMIT} reaches the wire
1.35\,ms before that response.
% src: facts-F6#43, facts-F6#45, facts-F6#46, facts-F6#49, facts-F6#50
This needs an unencrypted database link, so it runs on the plain capture
only.
% src: facts-F6#47

\emph{Break-glass is metadata only.} A size-only observer picks the largest
burst to the facility as the unlock, a rule validated on the plain capture.
% src: facts-F6#39
Under TLS it sees a 4356-byte unlock burst and no key or share: it learns
that, where and when a break-glass happened, and its size.
% src: facts-F6#40, facts-F6#41

\emph{Search leakage.} The run of Table~\ref{tab:wire} measured one search.
% src: facts-F6#53
The figures below come from two earlier runs, made before the 3-of-4
change (the search module is unchanged), each with 16 searches on one record,
one refused after expiry.
% src: facts-F6#91, facts-F6#57, facts-F6#59
On the plain wire all terms are visible, response size separates hit, miss
and refusal (43, 34, 95\,B), and the database link shows 15 lookups over 9
distinct tokens.
% src: facts-F6#60, facts-F6#61, facts-F6#62
Deterministic tokens link repeated searches within a session, the leakage of
deterministic searchable encryption~\cite{curtmola2006sse}, but terms searched
across expiry got new tokens.
% src: facts-F3#55, facts-F6#64, facts-F6#65
Under TLS no term or token is visible, yet response size still separates hit,
miss and refusal (60, 51, 112\,B) and request size is the term length plus
63\,B.
% src: facts-F6#68, facts-F6#69, facts-F6#71, facts-F6#72
A small sample on a different host showed no hit/miss timing difference
(66.1 vs 66.7\,ms median); there is no padding.
% src: facts-F6#74, facts-F6#75, facts-F6#77, facts-F6#92
